\begin{Definition}{环 \textmd{$R$} 上的多项式}\label{def:polynomial-over-ring}
  在环 $R$ 上, 关于变量 $x_{1},\ldots,x_{n}$ 的\textbf{多项式}是形如
  \[
    f=f(x_{1},\ldots,x_{n})=\sum_{\alpha\in\mathbb{N}^{n}}a_{\alpha}x^{\alpha},
  \]
  的表达式, 其中对每个 $\alpha\in\mathbb{N}^{n}$, 都有 $a_{\alpha}\in R$, 且除有限个 $\alpha$ 外, 均有 $a_{\alpha}=0$. 两个多项式
  \[
    f=\sum_{\alpha}a_{\alpha}x^{\alpha}
    \quad\text{和}\quad
    g=\sum_{\alpha}b_{\alpha}x^{\alpha}
  \]
  相等, 当且仅当对所有 $\alpha\in\mathbb{N}^{n}$, 都有 $a_{\alpha}=b_{\alpha}$.
\end{Definition}

\begin{Example}{多项式}\label{ex:polynomials-over-integers}
  \ExampleSideNote{下面是整数环 $\mathbb{Z}$ 上关于变量 $x$ 和 $y$ 的两个多项式:}%
    {多项式通常写成有限和, 省略所有系数为零的项.}
  \[
    f=xy^{2}+3xy+2
    \quad\text{和}\quad
    g=-xy+4.
  \]
  也可以把多项式 $f$ 和 $g$ 看作系数属于 $\mathbb{Q}$, $\mathbb{R}$, $\mathbb{C}$ 或任意其他包含 $\mathbb{Z}$ 的环的多项式.
  注意, 我们可以用熟悉的方式对 $f$ 和 $g$ 作加法和乘法, 从而得到新的多项式:
  \[
    \begin{gathered}
      f+g=xy^{2}+2xy+6,\\
      fg=-x^{2}y^{3}-3x^{2}y^{2}+4xy^{2}+10xy+8.
    \end{gathered}
  \]
\end{Example}

请读者在 Exercise~\ref{ex:polynomial-ring-laws} 中验证: 下一个定义中严格给出的加法和乘法, 使多项式的集合具有环结构.

\begin{Definition}{多项式环}\label{def:polynomial-ring}
  \textbf{多项式环} $R[x_{1},\ldots,x_{n}]$ 是环 $R$ 上关于变量 $x_{1},\ldots,x_{n}$ 的全体多项式的集合. 多项式加法和乘法定义如下:
  \[
    \left(\sum_{\alpha}a_{\alpha}x^{\alpha}\right)
    +\left(\sum_{\alpha}b_{\alpha}x^{\alpha}\right)
    =\sum_{\alpha}(a_{\alpha}+b_{\alpha})x^{\alpha}
  \]
  以及
  \[
    \left(\sum_{\alpha}a_{\alpha}x^{\alpha}\right)
    \left(\sum_{\alpha}b_{\alpha}x^{\alpha}\right)
    =\sum_{\alpha}\left(\sum_{\alpha_{1}+\alpha_{2}=\alpha}
      a_{\alpha_{1}}b_{\alpha_{2}}\right)x^{\alpha}.
  \]
\end{Definition}

加法单位元 $0\in R[x_{1},\ldots,x_{n}]$ 是对所有 $\alpha$ 均有 $a_{\alpha}=0$ 的多项式, 乘法单位元 $1\in R[x_{1},\ldots,x_{n}]$ 是满足下式的多项式:
\[
  a_{\alpha}=
  \begin{cases}
    1 & \text{若 }(\alpha_{1},\ldots,\alpha_{n})=(0,\ldots,0),\\
    0 & \text{若 }(\alpha_{1},\ldots,\alpha_{n})\ne(0,\ldots,0).
  \end{cases}
\]
